\documentclass{beamer}
\usepackage[utf8]{inputenc}

\usetheme{Madrid}
\usecolortheme{default}
\usepackage{amsmath,amssymb,amsfonts,amsthm}
\usepackage{txfonts}
\usepackage{tkz-euclide}
\usepackage{listings}
\usepackage{adjustbox}
\usepackage{array}
\usepackage{tabularx}
\usepackage{gvv}
\usepackage{lmodern}
\usepackage{circuitikz}
\usepackage{tikz}
\usepackage{graphicx}
\usepackage{mathtools}

\setbeamertemplate{page number in head/foot}[totalframenumber]

\usepackage{tcolorbox}
\tcbuselibrary{minted,breakable,xparse,skins}



\definecolor{bg}{gray}{0.95}
\DeclareTCBListing{mintedbox}{O{}m!O{}}{%
  breakable=true,
  listing engine=minted,
  listing only,
  minted language=#2,
  minted style=default,
  minted options={%
    linenos,
    gobble=0,
    breaklines=true,
    breakafter=,,
    fontsize=\small,
    numbersep=8pt,
    #1},
  boxsep=0pt,
  left skip=0pt,
  right skip=0pt,
  left=25pt,
  right=0pt,
  top=3pt,
  bottom=3pt,
  arc=5pt,
  leftrule=0pt,
  rightrule=0pt,
  bottomrule=2pt,
  toprule=2pt,
  colback=bg,
  colframe=orange!70,
  enhanced,
  overlay={%
    \begin{tcbclipinterior}
    \fill[orange!20!white] (frame.south west) rectangle ([xshift=20pt]frame.north west);
    \end{tcbclipinterior}},
  #3,
}
\lstset{
    language=C,
    basicstyle=\ttfamily\small,
    keywordstyle=\color{blue},
    stringstyle=\color{orange},
    commentstyle=\color{green!60!black},
    numbers=left,
    numberstyle=\tiny\color{gray},
    breaklines=true,
    showstringspaces=false,
}
\title{12.579}
\date{4th October, 2025}
\author{Puni Aditya - EE25BTECH11046}

\begin{document}

\frame{\titlepage}
\begin{frame}{Question}
Let V be the vector space of all 3 × 3 matrices with complex entries over the real
field. If
\begin{align*}
    W_1 &= \cbrak{\vec{A} \in V : \vec{A} = \bar{\vec{A}}^\top} \\
    W_2 &= \cbrak{\vec{A} \in V : \text{trace} \brak{\vec{A}} = 0}
\end{align*}
then the dimension of $W_1 + W_2$ is equal to \rule{2cm}{0.4pt}.
\end{frame}

\begin{frame}{Theoretical Solution}
The dimension of the sum of two subspaces is given by the formula:
\begin{align}
    \dim\brak{W_1 + W_2} = \dim\brak{W_1} + \dim\brak{W_2} - \dim\brak{W_1 \cap W_2}
\end{align}
The dimension of $V$, the space of $3 \times 3$ complex matrices over $\mathbb{R}$, is:
\begin{align}
    \dim\brak{V} = 3 \times 3 \times 2 = 18
\end{align}
\end{frame}

\begin{frame}{Theoretical Solution}
For $W_1$ (Hermitian matrices), the dimension is determined by 3 real diagonal elements and 3 complex upper-triangular elements.
\begin{align}
    \dim\brak{W_1} = 3\brak{1} + 3\brak{2} = 9
\end{align}
For $W_2$ (trace-zero matrices), the condition imposes two real constraints on the dimension of $V$.
\begin{align}
    \dim\brak{W_2} = \dim\brak{V} - 2 = 18 - 2 = 16
\end{align}
\end{frame}

\begin{frame}{Theoretical Solution}
For $W_1 \cap W_2$ (Hermitian, trace-zero matrices), the trace condition imposes one real constraint on the dimension of $W_1$.
\begin{align}
    \dim\brak{W_1 \cap W_2} = \dim\brak{W_1} - 1 = 9 - 1 = 8
\end{align}
Substituting these values into the dimension sum formula:
\begin{align}
    \dim\brak{W_1 + W_2} = 9 + 16 - 8 = 17
\end{align}
\end{frame}

\begin{frame}{Conclusion}
The dimension of $W_1 + W_2$ is 17.
\end{frame}

\end{document}
